% ============================================================
% Question Bank: ch09-04 Probability Models
% Topic Title: Probability Models
% CCSS: 7.SP.C.7
% Questions: 27 (15 MC + 6 SA + 6 Graphical)
% ============================================================

% --- Multiple Choice Questions (15) ---

\begin{practiceQuestion}{9-4-q01}{mc}
\begin{questionText}
A spinner has $5$ equal sections. Each section has the same chance of being selected. What type of probability model does this represent?
\end{questionText}
\choiceA{Non-uniform}
\choiceB{Uniform}
\choiceC{Experimental}
\choiceD{Biased}
\correctAnswer{B}
\explanation{A uniform probability model assigns equal probability to each outcome. Since all $5$ sections are equal, each has probability $\frac{1}{5}$, making this a uniform model.}
\end{practiceQuestion}

\begin{practiceQuestion}{9-4-q02}{mc}
\begin{questionText}
A grab bag has $6$ prizes worth $\$1$ each, $3$ prizes worth $\$5$ each, and $1$ grand prize worth $\$20$. What type of probability model describes drawing a prize by dollar value?
\end{questionText}
\choiceA{Uniform, because every prize has the same value}
\choiceB{Uniform, because each item has the same chance of being drawn}
\choiceC{Non-uniform, because the $\$1$ prizes are more likely to be drawn}
\choiceD{Non-uniform, because the $\$20$ prize is most likely}
\correctAnswer{C}
\explanation{There are different numbers of prizes at each value level. The probability of drawing a $\$1$ prize is $\frac{6}{10}$, a $\$5$ prize is $\frac{3}{10}$, and the $\$20$ prize is $\frac{1}{10}$. Because the probabilities are not equal, this is a non-uniform model.}
\end{practiceQuestion}

\begin{practiceQuestion}{9-4-q03}{mc}
\begin{questionText}
In any valid probability model, what must the probabilities of all outcomes add up to?
\end{questionText}
\choiceA{$0$}
\choiceB{$\dfrac{1}{2}$}
\choiceC{$1$}
\choiceD{It depends on the number of outcomes.}
\correctAnswer{C}
\explanation{A fundamental property of any probability model is that the sum of all outcome probabilities must equal $1$. This is true for both uniform and non-uniform models.}
\end{practiceQuestion}

\begin{practiceQuestion}{9-4-q04}{mc}
\begin{questionText}
A probability model has four outcomes with probabilities $0.35$, $0.25$, $0.20$, and $0.20$. Is this model valid?
\end{questionText}
\choiceA{No, because the probabilities are not all equal.}
\choiceB{No, because they do not add up to $1$.}
\choiceC{Yes, because each probability is between $0$ and $1$ and they sum to $1$.}
\choiceD{Yes, because there are exactly $4$ outcomes.}
\correctAnswer{C}
\explanation{Each probability is between $0$ and $1$, and $0.35 + 0.25 + 0.20 + 0.20 = 1.00$. Both conditions for a valid probability model are satisfied.}
\end{practiceQuestion}

\begin{practiceQuestion}{9-4-q05}{mc}
\begin{questionText}
Which set of probabilities does NOT represent a valid probability model for $3$ outcomes?
\end{questionText}
\choiceA{$\dfrac{1}{3},\; \dfrac{1}{3},\; \dfrac{1}{3}$}
\choiceB{$\dfrac{1}{2},\; \dfrac{1}{4},\; \dfrac{1}{4}$}
\choiceC{$\dfrac{1}{2},\; \dfrac{1}{3},\; \dfrac{1}{4}$}
\choiceD{$0.5,\; 0.3,\; 0.2$}
\correctAnswer{C}
\explanation{For a valid model, probabilities must sum to $1$. Check: $\frac{1}{2} + \frac{1}{3} + \frac{1}{4} = \frac{6}{12} + \frac{4}{12} + \frac{3}{12} = \frac{13}{12} > 1$. This exceeds $1$, so it is not valid.}
\end{practiceQuestion}

\begin{practiceQuestion}{9-4-q06}{mc}
\begin{questionText}
A teacher creates a probability model for a class game. The outcomes and their probabilities are: Win $= 0.2$, Tie $= 0.3$, and Lose $= ?$. What must $P(\text{Lose})$ equal?
\end{questionText}
\choiceA{$0.2$}
\choiceB{$0.5$}
\choiceC{$0.3$}
\choiceD{$0.7$}
\correctAnswer{B}
\explanation{Probabilities must sum to $1$. So $P(\text{Lose}) = 1 - 0.2 - 0.3 = 0.5$.}
\end{practiceQuestion}

\begin{practiceQuestion}{9-4-q07}{mc}
\begin{questionText}
A coin is weighted so that heads has a probability of $0.65$. After $200$ flips, heads appears $124$ times. How do the predicted and observed frequencies of heads compare?
\end{questionText}
\choiceA{Predicted: $130$; Observed: $124$. The observed is less.}
\choiceB{Predicted: $100$; Observed: $124$. The observed is greater.}
\choiceC{Predicted: $124$; Observed: $130$. The observed is greater.}
\choiceD{Predicted: $130$; Observed: $130$. They are equal.}
\correctAnswer{A}
\explanation{The predicted frequency is $0.65 \times 200 = 130$. The observed frequency is $124$. Since $124 < 130$, the observed is less than predicted — a normal small discrepancy.}
\end{practiceQuestion}

\begin{practiceQuestion}{9-4-q08}{mc}
\begin{questionText}
A standard number cube is a uniform probability model. A teacher rolls it $60$ times and gets: $1$: $12$, $2$: $11$, $3$: $9$, $4$: $8$, $5$: $10$, $6$: $10$. The predicted frequency for each number is $10$. Which statement is correct?
\end{questionText}
\choiceA{The cube must be unfair because the results are not all $10$.}
\choiceB{The small differences from $10$ are expected due to chance.}
\choiceC{The experiment should be discarded because it does not match.}
\choiceD{The model should be changed to non-uniform.}
\correctAnswer{B}
\explanation{In a uniform model, each outcome has the same probability but observed frequencies will vary slightly from predicted due to chance. The small deviations here are normal.}
\end{practiceQuestion}

\begin{practiceQuestion}{9-4-q09}{mc}
\begin{questionText}
A survey of $300$ students reveals that $150$ prefer pizza, $90$ prefer tacos, and $60$ prefer burgers for lunch. Which probability model best represents this data?
\end{questionText}
\choiceA{Uniform: $P(\text{pizza}) = P(\text{tacos}) = P(\text{burgers}) = \frac{1}{3}$}
\choiceB{Non-uniform: $P(\text{pizza}) = \frac{1}{2}$, $P(\text{tacos}) = \frac{3}{10}$, $P(\text{burgers}) = \frac{1}{5}$}
\choiceC{Non-uniform: $P(\text{pizza}) = \frac{1}{3}$, $P(\text{tacos}) = \frac{1}{3}$, $P(\text{burgers}) = \frac{1}{3}$}
\choiceD{Non-uniform: $P(\text{pizza}) = \frac{1}{2}$, $P(\text{tacos}) = \frac{1}{4}$, $P(\text{burgers}) = \frac{1}{4}$}
\correctAnswer{B}
\explanation{Based on the data: $P(\text{pizza}) = \frac{150}{300} = \frac{1}{2}$, $P(\text{tacos}) = \frac{90}{300} = \frac{3}{10}$, $P(\text{burgers}) = \frac{60}{300} = \frac{1}{5}$. These are not equal, so the model is non-uniform.}
\end{practiceQuestion}

\begin{practiceQuestion}{9-4-q10}{mc}
\begin{questionText}
A spinner is designed so that $P(\text{Red}) = 0.4$, $P(\text{Blue}) = 0.35$, and $P(\text{Green}) = 0.25$. If the spinner is spun $500$ times, what is the predicted frequency for Blue?
\end{questionText}
\choiceA{$125$}
\choiceB{$175$}
\choiceC{$200$}
\choiceD{$250$}
\correctAnswer{B}
\explanation{The predicted frequency is the probability times the number of trials: $0.35 \times 500 = 175$.}
\end{practiceQuestion}

\begin{practiceQuestion}{9-4-q11}{mc}
\begin{questionText}
Omar builds a probability model where the outcomes are A, B, C, and D. He assigns $P(\text{A}) = 0.1$, $P(\text{B}) = 0.4$, $P(\text{C}) = 0.3$. After $200$ trials, outcome D appears $38$ times. What is the predicted frequency of D?
\end{questionText}
\choiceA{$20$}
\choiceB{$30$}
\choiceC{$40$}
\choiceD{$50$}
\correctAnswer{C}
\explanation{First find $P(\text{D}) = 1 - 0.1 - 0.4 - 0.3 = 0.2$. The predicted frequency of D in $200$ trials is $0.2 \times 200 = 40$.}
\end{practiceQuestion}

\begin{practiceQuestion}{9-4-q12}{mc}
\begin{questionText}
A bag contains $2$ red, $5$ blue, and $3$ green tokens. A student says this is a uniform model. Which response is correct?
\end{questionText}
\choiceA{Correct — all tokens have the same chance of being drawn.}
\choiceB{Incorrect — the probabilities of the color outcomes are not equal.}
\choiceC{Correct — there are exactly $3$ outcomes.}
\choiceD{Incorrect — uniform models can only have $2$ outcomes.}
\correctAnswer{B}
\explanation{While each individual token is equally likely, the color outcomes have different probabilities: $P(\text{red}) = \frac{2}{10}$, $P(\text{blue}) = \frac{5}{10}$, $P(\text{green}) = \frac{3}{10}$. Because these are not equal, it is a non-uniform model for color.}
\end{practiceQuestion}

\begin{practiceQuestion}{9-4-q13}{mc}
\begin{questionText}
A game has $3$ outcomes. The developer wants outcome A to happen twice as often as outcome B, and outcome B to happen twice as often as outcome C. If $P(\text{C}) = \dfrac{1}{7}$, what is $P(\text{A})$?
\end{questionText}
\choiceA{$\dfrac{2}{7}$}
\choiceB{$\dfrac{3}{7}$}
\choiceC{$\dfrac{4}{7}$}
\choiceD{$\dfrac{1}{7}$}
\correctAnswer{C}
\explanation{If $P(\text{C}) = \frac{1}{7}$, then $P(\text{B}) = 2 \times \frac{1}{7} = \frac{2}{7}$ and $P(\text{A}) = 2 \times \frac{2}{7} = \frac{4}{7}$. Check: $\frac{4}{7} + \frac{2}{7} + \frac{1}{7} = 1$. \checkmark}
\end{practiceQuestion}

\begin{practiceQuestion}{9-4-q14}{mc}
\begin{questionText}
A candy machine dispenses colors with the following model: $P(\text{Red}) = 0.30$, $P(\text{Blue}) = 0.25$, $P(\text{Green}) = 0.25$, $P(\text{Yellow}) = 0.20$. After dispensing $400$ candies, $92$ were green. What is the discrepancy between the observed and predicted frequency for green?
\end{questionText}
\choiceA{$2$ fewer than predicted}
\choiceB{$8$ fewer than predicted}
\choiceC{$8$ more than predicted}
\choiceD{$2$ more than predicted}
\correctAnswer{B}
\explanation{Predicted green frequency: $0.25 \times 400 = 100$. Observed: $92$. The discrepancy is $100 - 92 = 8$, meaning green appeared $8$ fewer times than predicted.}
\end{practiceQuestion}

\begin{practiceQuestion}{9-4-q15}{mc}
\begin{questionText}
Which of the following best explains why observed frequencies often differ from predicted frequencies?
\end{questionText}
\choiceA{The probability model must be wrong.}
\choiceB{Random variation naturally causes differences in any finite experiment.}
\choiceC{Predicted frequencies are only guesses, not based on math.}
\choiceD{Observed frequencies are always more accurate than predicted ones.}
\correctAnswer{B}
\explanation{Random variation is expected in any experiment. With a finite number of trials, observed frequencies will not exactly match predicted ones, but they should be close. Larger sample sizes tend to reduce this discrepancy.}
\end{practiceQuestion}

% --- Short Answer Questions (6) ---

\begin{practiceQuestion}{9-4-q16}{sa}
\begin{questionText}
A probability model has $5$ equally likely outcomes labeled V, W, X, Y, and Z. What is the probability of each outcome? Express your answer as a simplified fraction.
\end{questionText}
\correctAnswer{$\frac{1}{5}$}
\explanation{In a uniform model, each outcome has the same probability. With $5$ outcomes, each has probability $\frac{1}{5}$.}
\end{practiceQuestion}

\begin{practiceQuestion}{9-4-q17}{sa}
\begin{questionText}
A probability model assigns $P(\text{A}) = 0.45$, $P(\text{B}) = 0.30$, and $P(\text{C}) = 0.25$. If the experiment is conducted $200$ times, what is the predicted frequency of outcome A?
\end{questionText}
\correctAnswer{$90$}
\explanation{Predicted frequency equals probability times the number of trials: $0.45 \times 200 = 90$.}
\end{practiceQuestion}

\begin{practiceQuestion}{9-4-q18}{sa}
\begin{questionText}
A model predicts that a spinner will land on red $\dfrac{2}{5}$ of the time. In $250$ spins, it lands on red $92$ times. What is the discrepancy between the observed and predicted frequency?
\end{questionText}
\correctAnswer{$8$ fewer}
\explanation{Predicted frequency: $\frac{2}{5} \times 250 = 100$. Observed: $92$. The discrepancy is $92 - 100 = -8$, meaning red appeared $8$ fewer times than predicted.}
\end{practiceQuestion}

\begin{practiceQuestion}{9-4-q19}{sa}
\begin{questionText}
A probability model for a board game has outcomes Win, Lose, and Draw with $P(\text{Win}) = 0.15$ and $P(\text{Draw}) = 0.25$. What is $P(\text{Lose})$?
\end{questionText}
\correctAnswer{$0.6$}
\explanation{All probabilities must sum to $1$. So $P(\text{Lose}) = 1 - 0.15 - 0.25 = 0.60$.}
\end{practiceQuestion}

\begin{practiceQuestion}{9-4-q20}{sa}
\begin{questionText}
A gumball machine has red, blue, and white gumballs. $P(\text{Red}) = \dfrac{3}{8}$ and $P(\text{Blue}) = \dfrac{3}{8}$. What is $P(\text{White})$? Express your answer as a simplified fraction.
\end{questionText}
\correctAnswer{$\frac{1}{4}$}
\explanation{$P(\text{White}) = 1 - \frac{3}{8} - \frac{3}{8} = 1 - \frac{6}{8} = \frac{2}{8} = \frac{1}{4}$.}
\end{practiceQuestion}

\begin{practiceQuestion}{9-4-q21}{sa}
\begin{questionText}
A fair die is rolled $120$ times. The number $3$ appears $18$ times. Is the observed frequency reasonably close to the predicted frequency? Answer Yes or No, and state the predicted frequency.
\end{questionText}
\correctAnswer{Yes; predicted is $20$}
\explanation{The predicted frequency is $\frac{1}{6} \times 120 = 20$. The observed frequency is $18$, which is very close to $20$. This small difference of $2$ is expected due to natural random variation.}
\end{practiceQuestion}

% --- Graphical Multiple Choice Questions (3) ---

\begin{practiceQuestion}{9-4-q22}{gmc}
\begin{questionText}
The spinner below shows a non-uniform probability model. What is $P(\text{Blue})$?

\begin{center}
\begin{tikzpicture}
  \fill[funRed!60] (0,0) -- (0:2.5) arc (0:180:2.5) -- cycle;
  \fill[funBlue!60] (0,0) -- (180:2.5) arc (180:270:2.5) -- cycle;
  \fill[funGreen!60] (0,0) -- (270:2.5) arc (270:360:2.5) -- cycle;
  \draw[thick] (0,0) circle (2.5);
  \draw[thick] (0,0) -- (0:2.5);
  \draw[thick] (0,0) -- (180:2.5);
  \draw[thick] (0,0) -- (270:2.5);
  \node at (90:1.4) {\textbf{Red}};
  \node at (90:0.8) {$\frac{1}{2}$};
  \node at (225:1.4) {\textbf{Blue}};
  \node at (225:0.8) {?};
  \node at (315:1.4) {\textbf{Green}};
  \node at (315:0.8) {$\frac{1}{4}$};
\end{tikzpicture}
\end{center}
\end{questionText}
\choiceA{$\dfrac{1}{2}$}
\choiceB{$\dfrac{1}{3}$}
\choiceC{$\dfrac{1}{4}$}
\choiceD{$\dfrac{1}{8}$}
\correctAnswer{C}
\explanation{Probabilities must sum to $1$. $P(\text{Blue}) = 1 - \frac{1}{2} - \frac{1}{4} = \frac{4}{4} - \frac{2}{4} - \frac{1}{4} = \frac{1}{4}$.}
\end{practiceQuestion}

\begin{practiceQuestion}{9-4-q23}{gmc}
\begin{questionText}
The double bar graph compares predicted and observed frequencies for outcomes of a game played $100$ times.

\begin{center}
\begin{tikzpicture}
  \begin{axis}[
    ybar, bar width=12pt, ymin=0, ymax=55,
    ylabel={Frequency},
    symbolic x coords={Win, Lose, Draw},
    xtick=data,
    width=8cm, height=6cm,
    enlarge x limits=0.35,
    legend style={at={(0.5,1.05)}, anchor=south, legend columns=2},
    nodes near coords
  ]
  \addplot[fill=funBlue!80] coordinates {(Win,22) (Lose,48) (Draw,30)};
  \addplot[fill=funBlue!25] coordinates {(Win,20) (Lose,50) (Draw,30)};
  \legend{Observed, Predicted}
  \end{axis}
\end{tikzpicture}
\end{center}

Which outcome has the largest discrepancy between observed and predicted?
\end{questionText}
\choiceA{Win only}
\choiceB{Lose only}
\choiceC{Draw only}
\choiceD{Win and Lose — both differ by $2$}
\correctAnswer{D}
\explanation{Win: observed $22$, predicted $20$ (difference of $2$). Lose: observed $48$, predicted $50$ (difference of $2$). Draw: observed $30$, predicted $30$ (difference of $0$). Win and Lose both have the largest discrepancy of $2$.}
\end{practiceQuestion}

\begin{practiceQuestion}{9-4-q24}{gmc}
\begin{questionText}
The table below shows a probability model for a mystery bag of colored tiles.

\begin{center}
\begin{tabular}{|c|c|c|c|c|}
\hline
\textbf{Color} & Red & Blue & Green & Yellow \\
\hline
\textbf{Probability} & $0.40$ & $0.25$ & $0.20$ & $0.15$ \\
\hline
\end{tabular}
\end{center}

If $80$ tiles are drawn (with replacement), what is the predicted number of yellow tiles?
\end{questionText}
\choiceA{$10$}
\choiceB{$12$}
\choiceC{$15$}
\choiceD{$20$}
\correctAnswer{B}
\explanation{Predicted frequency = probability $\times$ trials = $0.15 \times 80 = 12$ yellow tiles.}
\end{practiceQuestion}

% --- Graphical Short Answer Questions (3) ---

\begin{practiceQuestion}{9-4-q25}{gsa}
\begin{questionText}
A spinner is divided into sections as shown. The probability of each color is given.

\begin{center}
\begin{tikzpicture}
  \fill[funRed!60] (0,0) -- (0:2.5) arc (0:144:2.5) -- cycle;
  \fill[funBlue!60] (0,0) -- (144:2.5) arc (144:252:2.5) -- cycle;
  \fill[funGreen!60] (0,0) -- (252:2.5) arc (252:360:2.5) -- cycle;
  \draw[thick] (0,0) circle (2.5);
  \draw[thick] (0,0) -- (0:2.5);
  \draw[thick] (0,0) -- (144:2.5);
  \draw[thick] (0,0) -- (252:2.5);
  \node at (72:1.4) {\textbf{Red}};
  \node at (72:0.8) {$P = 0.40$};
  \node at (198:1.4) {\textbf{Blue}};
  \node at (198:0.8) {$P = 0.30$};
  \node at (306:1.4) {\textbf{Green}};
  \node at (306:0.8) {$P = 0.30$};
\end{tikzpicture}
\end{center}

If this spinner is spun $150$ times, what is the predicted frequency of landing on red?
\end{questionText}
\correctAnswer{$60$}
\explanation{Predicted frequency = probability $\times$ trials = $0.40 \times 150 = 60$.}
\end{practiceQuestion}

\begin{practiceQuestion}{9-4-q26}{gsa}
\begin{questionText}
The table below shows the results of a survey asking $500$ commuters how they get to work.

\begin{center}
\begin{tabular}{|l|c|c|}
\hline
\textbf{Mode} & \textbf{Predicted (\%)} & \textbf{Observed Count} \\
\hline
Drive & $50\%$ & $262$ \\
\hline
Bus & $25\%$ & $118$ \\
\hline
Walk & $15\%$ & $80$ \\
\hline
Bike & $10\%$ & $40$ \\
\hline
\end{tabular}
\end{center}

Which mode of transportation has the greatest discrepancy between predicted and observed count? State the difference.
\end{questionText}
\correctAnswer{Drive; $12$ more}
\explanation{Predicted: Drive $= 250$, Bus $= 125$, Walk $= 75$, Bike $= 50$. Observed: Drive $= 262$ (difference of $12$), Bus $= 118$ (difference of $7$), Walk $= 80$ (difference of $5$), Bike $= 40$ (difference of $10$). Drive has the greatest discrepancy of $12$.}
\end{practiceQuestion}

\begin{practiceQuestion}{9-4-q27}{gsa}
\begin{questionText}
The bar graph shows the results of rolling a loaded die $240$ times.

\begin{center}
\begin{tikzpicture}
  \begin{axis}[
    ybar, bar width=16pt, ymin=0, ymax=70,
    ylabel={Frequency},
    xlabel={Outcome},
    symbolic x coords={1, 2, 3, 4, 5, 6},
    xtick=data,
    nodes near coords,
    width=9cm, height=6cm,
    enlarge x limits=0.15
  ]
  \addplot[fill=funBlue!70] coordinates {(1,30) (2,32) (3,28) (4,35) (5,55) (6,60)};
  \end{axis}
\end{tikzpicture}
\end{center}

Based on these results, develop a probability model by finding $P(5)$ and $P(6)$. Express each as a simplified fraction.
\end{questionText}
\correctAnswer{$P(5) = \frac{11}{48}$; $P(6) = \frac{1}{4}$}
\explanation{$P(5) = \frac{55}{240} = \frac{11}{48}$ and $P(6) = \frac{60}{240} = \frac{1}{4}$. This die is non-uniform — outcomes $5$ and $6$ appear more often than a fair die would predict ($\frac{1}{6} \approx 40$ each).}
\end{practiceQuestion}
